\section{Pruebas}

\subsection{Pruebas automáticas}

\begin{tabularx}{\linewidth}{@{}L L L@{}}
\toprule
Suite & Alcance & Resultado \\
\midrule
Tests unitarios de la app & Reglas de puntos, casos de uso, repositorios, workers, ViewModels y escenarios de uso real & 207 tests. Corren en la JVM, sin emulador \\
Tests del backend & Autenticación, catálogo, canjes, uso con PIN, libro de puntos y concurrencia & 78 casos declarados \\
Cobertura de \texttt{domain} & Umbral de 80 \% verificado por Kover en el CI & 100 \% de las líneas (medido con Kover el 19/11/2026). El CI falla si baja del 80 \% \\
Tests instrumentados & Room, workers con work-testing y RNF04 & 17 de 17 aprobados en el emulador (Android 16, API 36) el 20/11/2026. En API 28 no se ejecutaron \\
Tests de UI & Pantallas en cada estado y gesto de confirmación & No implementados: las pantallas se recorrieron a mano en el emulador (plan de pruebas manuales) \\
\bottomrule
\end{tabularx}

Los tests unitarios de la app y los del backend corren en GitHub Actions en cada pull request y en cada
cambio que llega a \texttt{dev}.

\subsection{Escenarios de uso real}
\label{sec:escenarios}

Cada regla de puntos se prueba con la situación que la motiva. El test recorre el camino completo: los
eventos que registraría Android, las pausas que arma la app y los puntos y el motivo que asignan las
reglas (\texttt{ScenariosTest}, 16 casos que corren en el CI).

\begin{table}[H]
\centering
\begin{tabularx}{\linewidth}{@{}L L L@{}}
\toprule
Situación & Regla que la cubre & Resultado \\
\midrule
Deja el teléfono en la mesa para cenar una hora & 1 punto por minuto válido & +60 \\
Una notificación prende la pantalla sin desbloquear & La pausa sigue hasta desbloquear & Una sola pausa, +60 \\
Mira la hora 20 segundos y lo vuelve a bloquear & Una pausa de menos de 1 min no se guarda & No aparece en el historial \\
Lo deja 20 minutos para hacer un mandado & Mínimo de 30 min los días hábiles & 0, muy corta \\
Duerme una siesta de 40 minutos un sábado & Mínimo de 45 min el fin de semana & 0, muy corta \\
Duerme toda la noche & De 23 a 9 h no suma & 0, horario de descanso \\
Lee de 22:30 a 23:45 & Solo cuentan los minutos antes de las 23 & +30 \\
Se olvida el teléfono en un cajón 5 horas & Tope de 3 h por pausa & +180 \\
Hace tres pausas de 2 horas el mismo día & Tope de 240 puntos por día & +120, +120 y 0 \\
Pone música y bloquea el teléfono & No suma si un servicio reciente de otra app estuvo activo la mitad de la pausa o más & 0, había música \\
Atiende una llamada larga con la pantalla apagada & La misma regla & 0, había una llamada \\
La música se corta a los 10 minutos & El servicio no llega a la mitad de la pausa & +60 \\
Tiene una VPN o un contador de pasos prendido todo el día & Se ignoran los servicios que arrancaron más de 30 min antes de bloquear & +60 \\
El teléfono se apaga por batería a la mitad & El tiempo apagado no cuenta & +40 \\
Cambia la hora del teléfono para sumar más & Una pausa con cambio de hora no suma & 0, cambió la hora \\
La app se cierra justo después de guardar & Releer una pausa no la duplica ni le baja los puntos (RNF04) & Mismo resultado \\
\bottomrule
\end{tabularx}
\caption{Escenarios de uso real y resultado de las reglas.}
\end{table}

Estos tests prueban la lógica de la app a partir de los eventos. Que Android entregue esos eventos y
cuándo corre el worker con las restricciones de batería se comprueba en el emulador, con el plan de
pruebas manuales.

\subsection{Plan de pruebas manuales}

El plan sigue la app desde que se instala hasta el canje. Se ejecutó el 20/11/2026 en el emulador de
Android Studio (Pixel 8, Android 16, API 36), porque ningún integrante tiene un teléfono Android, con el
backend local y la cuenta de una persona nueva. Cada caso se marca como aprobado, fallado o no ejecutado.

\begin{enumerate}
\item Instalar la app y completar el primer uso.
\item Entrar con email. Con Google, si el cliente ID está configurado.
\item Dar el acceso a uso y comprobar que la app detecta la vuelta desde Ajustes.
\item Bloquear el teléfono durante más de 30 minutos un día hábil y comprobar que la pausa suma.
\item Comprobar que una pausa dentro del horario de descanso no suma y que su motivo aparece en Actividad.
\item Canjear un premio con saldo suficiente y confirmar la entrega con el PIN del local.
\item Intentar usar el mismo código otra vez y comprobar que no pasa nada.
\item Intentar canjear en modo avión y comprobar que no se descuentan puntos.
\item Dejar pasar varios días sin abrir la app y comprobar que las pausas se acreditan al abrirla.
\item Rechazar los permisos y comprobar que la app sigue funcionando y lo indica.
\end{enumerate}

\begin{table}[H]
\centering
\begin{tabularx}{\linewidth}{@{}c L L@{}}
\toprule
Caso & Resultado & Observaciones \\
\midrule
1 & Aprobado & Splash, onboarding de tres pasos y acceso. La segunda vez el onboarding no aparece. En el paso 2 las tarjetas salían recortadas: se corrigió \\
2 & Aprobado & Cuenta nueva con email. Google no se probó porque el cliente ID no está configurado \\
3 & Aprobado & Sin el acceso a uso aparece el frame 07; al darlo desde Ajustes, la app sigue sola a Inicio \\
4 & Aprobado & Una pausa simulada de 60 min sumó 60 puntos, llegó el aviso ``Break validado'' y se vio en Inicio y Actividad. Pausas reales de 70 s se detectaron y quedaron como ``muy corta'' \\
5 & No ejecutado & Cubierto por los tests de escenarios (``duerme toda la noche'') \\
6 & Aprobado & El código lo generó el servidor, con cuenta regresiva y hora en vivo. Con un PIN incorrecto quedan 4 intentos; con el correcto, la entrega se confirma \\
7 & No ejecutado & Cubierto por los tests del backend (dos usos simultáneos, solo uno pasa) \\
8 & Aprobado & En modo avión aparece el aviso sin conexión y el canje dice que hace falta red, sin descontar puntos \\
9 & No ejecutado & Requiere días reales sin abrir la app; la retención de eventos de Android lo cubre en teoría \\
10 & Aprobado & Rechazar el acceso a uso lleva al frame 07; ``Ahora no'' deja ver premios. La exclusión de batería abre el diálogo del sistema y queda activa \\
\bottomrule
\end{tabularx}
\caption{Resultados del plan de pruebas manuales en el emulador.}
\end{table}

El recorrido encontró y corrigió tres errores: con el bloqueo de pantalla en ``Ninguno'' Android no emite
los eventos de bloqueo y no se detectaba ninguna pausa (ahora se usa el apagado y el encendido de la
pantalla); la ubicación nunca llegaba en el emulador y Explorar no se ordenaba por cercanía (ahora usa la
última conocida tras 5 s); y la racha de Inicio decía ``1 días''.

\subsection{Batería, arranque y privacidad}

El consumo de batería en 24 horas (RNF01) y el arranque en frío con Macrobenchmark (RNF07) no se
midieron: en el emulador la batería es simulada y los tiempos no representan un teléfono real. Quedan como
trabajo futuro. Por diseño, la app no usa servicios en primer plano ni ubicación en segundo plano, y lee
los eventos solo cuando corre el worker o se abre la app.

La revisión de privacidad y seguridad se hizo sobre el código el 20/11/2026:

\begin{itemize}
\item El build de release no permite tráfico sin cifrar; solo el de debug habla con el backend local.
\item El código no escribe logs y el token de sesión se guarda cifrado con una clave del Keystore.
\item Los permisos son los mínimos: datos de uso, notificaciones y ubicación aproximada en primer plano.
No se declara ubicación en segundo plano ni servicios en primer plano.
\item Al servidor solo llegan los totales de puntos por día, sin pausas ni horarios.
\item Se encontró que la copia de seguridad automática de Android subía la base con las pausas a la nube
del usuario. Se desactivó (\texttt{allowBackup="false"}).
\end{itemize}
